% Syntra LaTeX Template
% For Arabic support, use XeLaTeX with:
%   \usepackage{fontspec}
%   \usepackage{polyglossia}
%   \setdefaultlanguage{english}
%   \setotherlanguage{arabic}
%   \newfontfamily\arabicfont[Script=Arabic]{Amiri}

\documentclass[12pt,a4paper]{article}

% Font and encoding
\usepackage[utf8]{inputenc}
\usepackage[T1]{fontenc}

% Math packages
\usepackage{amsmath,amssymb}

% Table support
\usepackage{longtable,booktabs}

% Graphics and links
\usepackage{graphicx}
\usepackage{hyperref}
\usepackage{geometry}

% Page geometry
\geometry{
    left=2.5cm,
    right=2.5cm,
    top=3cm,
    bottom=3cm
}

% Fix pandoc tightlist issue
\providecommand{\tightlist}{\setlength{\itemsep}{0pt}\setlength{\parskip}{0pt}}

% Hyperref settings
\hypersetup{
    colorlinks=true,
    linkcolor=blue,
    filecolor=magenta,
    urlcolor=cyan,
    pdftitle={Syntra Research Documentation},
    pdfauthor={Syntra-Env},
    pdfpagemode=FullScreen
}

\begin{document}

-\/-\/- title: UOR v3.5.0 (The Prime Framework) prev\_title: Goals
prev\_url: 02-goals.html next\_title: Alternatives next\_url:
04-alternatives.html -\/-\/- Syntra aligns with \textbf{UOR v3.5.0},
implementing the \textbf{Prime Framework} for universal number
representation. This ensures \textbf{Observer Invariance} and
\textbf{Geometric Consistency} by encoding all information through
complete prime factorization.

\subsection{Universal Coordinate
Projection}\label{universal-coordinate-projection}

Content is projected onto the triadic UOR manifold using
prime-factorized Universal Numbers. Similarity is measured via
\textbf{Geometric Congruence} across prime spectra.

\begin{Shaded}
\begin{Highlighting}[]
\NormalTok{c1 }\OperatorTok{=}\NormalTok{ uor.triality.project\_coordinate(}\StringTok{"Hello world"}\NormalTok{)}
\NormalTok{c2 }\OperatorTok{=}\NormalTok{ uor.triality.project\_coordinate(}\StringTok{"Goodbye world"}\NormalTok{)}
\NormalTok{Coordinates are now stored }\ImportTok{as}\NormalTok{ \{prime: exponent\} factorizations}
\BuiltInTok{print}\NormalTok{(}\StringTok{"Coord 1 (Extension):"}\NormalTok{, c1[}\StringTok{\textquotesingle{}extension\textquotesingle{}}\NormalTok{])}
\BuiltInTok{print}\NormalTok{(}\StringTok{"Coord 2 (Extension):"}\NormalTok{, c2[}\StringTok{\textquotesingle{}extension\textquotesingle{}}\NormalTok{])}
\BuiltInTok{print}\NormalTok{()}
\NormalTok{resonance }\OperatorTok{=}\NormalTok{ uor.triality.compute\_resonance(c1, c2)}
\BuiltInTok{print}\NormalTok{(}\SpecialStringTok{f"Universal Resonance Fidelity: }\SpecialCharTok{\{}\NormalTok{resonance}\SpecialCharTok{:.4f\}}\SpecialStringTok{"}\NormalTok{)}
\end{Highlighting}
\end{Shaded}

\subsection{Holonomic Field --- Prime
Congruence}\label{holonomic-field-prime-congruence}

The \textbf{Holonomic Field} utilizes prime factorization to compute
observer-invariant congruence. \textbf{Semantic Tension} measures local
curvature within the prime-factorized semantic lattice.

\begin{Shaded}
\begin{Highlighting}[]
\NormalTok{dataset\_raw }\OperatorTok{=}\NormalTok{ [}
  \StringTok{"artificial intelligence"}\NormalTok{, }\StringTok{"machine learning"}\NormalTok{,}
  \StringTok{"deep neural network"}\NormalTok{, }\StringTok{"cooking recipe"}
\NormalTok{]}
\NormalTok{manifold\_points }\OperatorTok{=}\NormalTok{ [\{ }\StringTok{\textquotesingle{}label\textquotesingle{}}\NormalTok{: s, }\StringTok{\textquotesingle{}coordinate\textquotesingle{}}\NormalTok{: uor.triality.project\_coordinate(s) \} }\ControlFlowTok{for}\NormalTok{ s }\KeywordTok{in}\NormalTok{ dataset\_raw]}
\NormalTok{query }\OperatorTok{=}\NormalTok{ uor.triality.project\_coordinate(}\StringTok{"neural network training"}\NormalTok{)}
\NormalTok{neighbors }\OperatorTok{=}\NormalTok{ uor.field.manifold\_proximity(query, manifold\_points, }\DecValTok{3}\NormalTok{)}
\BuiltInTok{print}\NormalTok{(}\StringTok{"Manifold Proximity (Prime{-}Aware):"}\NormalTok{)}
\ControlFlowTok{for}\NormalTok{ n }\KeywordTok{in}\NormalTok{ neighbors:}
    \BuiltInTok{print}\NormalTok{(}\SpecialStringTok{f"  }\SpecialCharTok{\{}\NormalTok{n[}\StringTok{\textquotesingle{}item\textquotesingle{}}\NormalTok{][}\StringTok{\textquotesingle{}label\textquotesingle{}}\NormalTok{]}\SpecialCharTok{\}}\SpecialStringTok{ {-}\textgreater{} }\SpecialCharTok{\{}\NormalTok{n[}\StringTok{\textquotesingle{}score\textquotesingle{}}\NormalTok{]}\SpecialCharTok{:.3f\}}\SpecialStringTok{"}\NormalTok{)}
\end{Highlighting}
\end{Shaded}

\subsection{Gauge Stability \& Topology}\label{gauge-stability-topology}

The topology module assesses \textbf{Gauge Stability} by measuring
resonance fidelity across edges, ensuring mathematical consistency
across the manifold.

\begin{Shaded}
\begin{Highlighting}[]
\NormalTok{topology }\OperatorTok{=}\NormalTok{ uor.topology.build\_topology(}
\NormalTok{  [}
\NormalTok{    \{ }\StringTok{\textquotesingle{}id\textquotesingle{}}\NormalTok{: }\StringTok{"n1"}\NormalTok{, }\StringTok{\textquotesingle{}label\textquotesingle{}}\NormalTok{: }\StringTok{"Source A"}\NormalTok{ \},}
\NormalTok{    \{ }\StringTok{\textquotesingle{}id\textquotesingle{}}\NormalTok{: }\StringTok{"n2"}\NormalTok{, }\StringTok{\textquotesingle{}label\textquotesingle{}}\NormalTok{: }\StringTok{"Target B"}\NormalTok{ \}}
\NormalTok{  ],}
\NormalTok{  [}
\NormalTok{    \{ }\StringTok{\textquotesingle{}from\textquotesingle{}}\NormalTok{: }\StringTok{"n1"}\NormalTok{, }\StringTok{\textquotesingle{}to\textquotesingle{}}\NormalTok{: }\StringTok{"n2"}\NormalTok{, }\StringTok{\textquotesingle{}relation\textquotesingle{}}\NormalTok{: }\StringTok{"functor"}\NormalTok{ \}}
\NormalTok{  ]}
\NormalTok{)}
\NormalTok{stability }\OperatorTok{=}\NormalTok{ uor.topology.assess\_gauge\_stability(topology)}
\BuiltInTok{print}\NormalTok{(}\SpecialStringTok{f"Topology Gauge Stability: }\SpecialCharTok{\{}\NormalTok{stability[}\StringTok{\textquotesingle{}resonance\_fidelity\textquotesingle{}}\NormalTok{]}\SpecialCharTok{:.4f\}}\SpecialStringTok{"}\NormalTok{)}
\end{Highlighting}
\end{Shaded}


\end{document}
